\documentclass[10pt]{beamer}

\usetheme{metropolis}
\definecolor{WiLabRed}{RGB}{197,18,48}
\setbeamercolor{frametitle}{fg=white,bg=WiLabRed}
\setbeamercolor{progress bar}{fg=WiLabRed!90}
\setbeamercolor{title separator}{fg=WiLabRed!90}
\setbeamercolor{progress bar in section page}{fg=WiLabRed!90}
\setbeamercolor{background canvas}{bg=white}
\setbeamercolor{alerted text}{fg=WiLabRed!90}

\usepackage{appendixnumberbeamer}

\usepackage{booktabs}
\usepackage[scale=2]{ccicons}

\usepackage{pgfplots}
\usepgfplotslibrary{dateplot}

\usepackage{xspace}
\newcommand{\themename}{\textbf{\textsc{metropolis}}\xspace}

\usepackage{marvosym}
%\usepackage{subfig}
\usepackage{graphicx}\graphicspath{{images/}}
\usepackage{subcaption}
\usepackage[framed]{./support/mcode}
\usepackage{listings}

\title{Lecture 05: PSD \& LTI Systems}
\subtitle{\textit{Software Defined Radio for Engineers} (Collins~\textit{et al.}), \textsection{3.6}}
\date{}
\author{\textbf{Alexander M. Wyglinski, Ph.D.}}
\institute{ \vspace*{1in}\hfill\includegraphics[height=1.125cm]{wilab_logo-A70916.eps} \qquad \includegraphics[height=1.125cm]{WPI_Inst_Prim_FulClr.eps}}
% \titlegraphic{\hfill\includegraphics[height=1.5cm]{logo.pdf}}

% Foot for all slides
\setbeamertemplate{frame footer}{\tiny \copyright~2018 by Alexander Wyglinski. This work is licensed under the Creative Commons Attribution-ShareAlike 4.0 International License. To view a copy of this license, visit http://creativecommons.org/licenses/by-sa/4.0/.}

\begin{document}

%\captionsetup[subfigure]{labelformat=empty}

%%%%%%%%%%%%%%%%%%%%%%%%%%%%%%%%%%%%%%%%%%%%%%%%%%%%%%%%%%

\maketitle




%%%%%%%%%%%%%%%%%%%%%%%%%%%%%%%%%%%%%%%%%%%%%%%%%%%%%%%%%%

\begin{frame}[fragile]{Frequency Domain Perspective}

    \begin{itemize}
        \item It is sometimes more convenient to study a communication system in the frequency domain rather than the time domain
        \begin{itemize}
            \item Mathematical analysis is more tractable
            \item Operations such as convolution are transformed into simple multiplications
        \end{itemize}
        \item Physically it makes sense to study wireless transmissions in the frequency domain
        \begin{itemize}
            \item Digital communications is the transferral of data based on changes of electromagnetic wave characteristics such as frequency, amplitude, and phase
        \end{itemize}
    \end{itemize}

\end{frame}

\begin{frame}[fragile]{Fourier Transform}

    \begin{itemize}
        \item Mathematical relationship between a time domain waveform and weighted sum of sinusoidal components that constitute it, i.e., frequency domain representation
        \item Translating between time and frequency domains is achieved using the Fourier transform and inverse Fourier transform:
            \begin{equation}
                  H(f)=\int_{-\infty}^{\infty}h(t)e^{-j2\pi{f}{t}}dt
            \end{equation}
            \begin{equation}
                  h(t)=\int_{-\infty}^{\infty}H(f)e^{j2\pi{f}{t}}df
            \end{equation}
    \end{itemize}

\end{frame}

\begin{frame}[fragile]{Einstein-Wiener-Khinchin Theorem}

    \begin{itemize}
        \item In particular, we are interested in the {\it power spectral density} (PSD) of a signal, which is related to the autocorrelation function via the {\it Einstein-Wiener-Khintchine} (EWK) Relations:
            \begin{equation}
                  S_X(f)=\int_{-\infty}^{\infty}R_X(\tau)e^{-j2\pi{f}{\tau}}d\tau
            \end{equation}
            \begin{equation}                 
                  R_X(\tau)=\int_{-\infty}^{\infty}S_X(f)e^{j2\pi{f}{\tau}}df
            \end{equation} 
        \item Relating the PSD between input $x(t)$ and output $y(t)$ of a system $h(t)$, we have the following very important result:
            \begin{equation}
                  S_Y(f)=|H(f)|^2S_X(f)
            \end{equation}                
    \end{itemize}
\end{frame}

\begin{frame}[fragile]{Impulse Train Model}

    \begin{itemize}
        \item Signal analysis and knowledge of the input PSD can help determine the output PSD of the transmitter 
        \begin{itemize}
            \item Determine the autocorrelation function of $B_n$ from $I_n$ and then apply EWK relations
        \end{itemize}
    \end{itemize}

\end{frame}

\begin{frame}[fragile]{PSD Properties}
  
  \begin{itemize}
    \item Zero frequency scenario:
    \begin{equation}
        S_X(0)=\int_{-\infty}^{\infty}R_X(\tau)d\tau\nonumber
    \end{equation}
    \item Mean-squared value:
    \begin{equation}
        E\{x^2(t)\}=\int_{-\infty}^{\infty}S_X(f)df\nonumber
    \end{equation}
    \item Non-negative PSD: $S_X(f)\ge{0},~\forall{f}$
    \item Real-valued process: $S_X(-f)=S_X(f)$
    \item Normalized PSD associated with a PDF:
    \begin{equation}
        p_X(f)=\frac{S_X(f)}{\int_{-\infty}^{\infty}S_X(f)df}\nonumber
    \end{equation}
  \end{itemize}
  
\end{frame}

\begin{frame}[fragile]{PSD Example}

  \begin{itemize}
    \item Find $S_X(f)$ of the following random process:
    \begin{equation}
        x(t)=A\cos(2\pi{f_c}t+\Theta),\nonumber
    \end{equation}
    where $\Theta$ is uniformly distributed over the interval $[-\pi,\pi]$.
    \item First solve for $R_X(\tau)$ by using:
    \begin{equation}
        R_X(\tau)=E\{x(t+\tau)x(t)\}=\frac{1}{2}A^2\cos(2\pi{f_c}\tau)\nonumber
    \end{equation}
    \item Then solve for the PSD using EWK relations:
    \begin{equation}
    \begin{split}
        S_X(f)&=\int_{-\infty}^{\infty}R_X(\tau)e^{-j2\pi{f}\tau}d\tau=\int_{-\infty}^{\infty}\frac{1}{2}A^2\cos(2\pi{f_c}\tau)e^{-j2\pi{f}\tau}d\tau\\
        &=\int_{-\infty}^{\infty}\frac{1}{4}A^2\left(e^{j2\pi{f_c}\tau}+e^{-j2\pi{f_c}\tau}\right)e^{-j2\pi{f}\tau}d\tau\\
        &=\frac{A^2}{4}(\delta(f-f_c)+\delta(f+f_c))\nonumber
    \end{split}
    \end{equation}
  \end{itemize}
  
\end{frame}

\begin{frame}[fragile]{Identifying PSD Features}

    \begin{itemize}
        \item PSD characteristics can reveal a substantial amount of info regarding a signal
        \begin{itemize}
            \item Modulation format
            \item Pulse shape filtering
            \item Transmission bandwidth
        \end{itemize}
        \item Time-varying behavior indicates network traffic levels
        \item PSD shape can uniquely identify specific wireless standards
    \end{itemize}

\end{frame}






\end{document}
